\chapter{Clases, Objetos y Atributos de Instancia}

\section{Definición de una Clase en Python}
Para declarar una clase en Python se utiliza la palabra reservada \texttt{class}, seguida del nombre de la clase
(escrito convencionalmente en formato \textit{CamelCase}) y dos puntos.
Si deseamos definir una clase vacía sin atributos ni métodos iniciales, empleamos la sentencia \texttt{pass} como
marcador de posición sintáctico.

\begin{lstlisting}[caption={Definición de una clase vacía y su instanciación.}, label={list:clase_vacia}]
class Vehicle:
    """Clase base que representa un vehiculo generico."""
    pass

# Inspeccion del objeto clase
print(Vehicle)
\end{lstlisting}

En el \sect{list:clase_vacia} se observa el uso de un \textit{docstring} (cadena literal encerrada entre tres comillas),
el cual permite documentar la función o propósito de la clase directamente en el código fuente.

\section{El Constructor \texttt{\_\_init\_\_} y el Parámetro \texttt{self}}
Cuando se crea una nueva instancia de una clase, Python ejecuta automáticamente un método especial de inicialización
llamado \texttt{\_\_init\_\_}. Este método actúa como el constructor de la clase, reservando espacio y asignando los
valores iniciales a los atributos del objeto.

El primer parámetro de todo método de instancia en Python debe ser obligatoriamente \texttt{self}. El parámetro
\texttt{self} es una referencia explícita a la instancia concreta sobre la cual se está ejecutando el método.

\begin{lstlisting}[caption={Uso del constructor y asignación de atributos de instancia.}, label={list:init_atributos}]
class Vehicle:
    """Representa un vehiculo con velocidad maxima y kilometraje."""
    
    def __init__(self, name, max_speed, mileage):
        self.name = name
        self.max_speed = max_speed
        self.mileage = mileage

# Instanciacion de objetos independientes
v1 = Vehicle("Tesla Model S", 250, 18)
v2 = Vehicle("Toyota Corolla", 180, 45000)

print(f"Vehiculo 1: {v1.name}, Maxima: {v1.max_speed} km/h")
print(f"Vehiculo 2: {v2.name}, Maxima: {v2.max_speed} km/h")
\end{lstlisting}

Al ejecutar la línea \texttt{v1 = Vehicle("Tesla Model S", 250, 18)}, Python invoca internamente a
\texttt{Vehicle.\_\_init\_\_(v1, "Tesla Model S", 250, 18)}. La notación de punto (\texttt{v1.name})
permite acceder a los atributos específicos almacenados en esa instancia.

\section{Métodos de Instancia}
Los métodos de instancia son funciones definidas dentro de una clase que operan sobre los datos almacenados
en los atributos de la instancia actual. A través del parámetro \texttt{self}, los métodos pueden leer o modificar
el estado interno del objeto.

\subsection{Cálculos Derivados a partir de Atributos}
Un patrón habitual en POO es la creación de métodos que calculan nueva información utilizando los atributos
almacenados en la instancia.

\begin{lstlisting}[caption={Clase Rectangle con métodos de cálculo geométrico.}, label={list:rectangle}]
class Rectangle:
    def __init__(self, length, width):
        self.length = length
        self.width = width

    def area(self):
        """Calcula el area del rectangulo."""
        return self.length * self.width

    def perimeter(self):
        """Calcula el perimetro del rectangulo."""
        return 2 * (self.length + self.width)

rect = Rectangle(10, 4)
print(f"Area: {rect.area()}")
print(f"Perimetro: {rect.perimeter()}")
\end{lstlisting}

Al llamar a \texttt{rect.area()}, no es necesario pasar el argumento \texttt{self} explícitamente, ya que Python
asocia automáticamente el objeto \texttt{rect} invocante.

\subsection{Validación de Datos y Reglas de Negocio}
Los métodos también se utilizan para imponer restricciones y validar las operaciones sobre el estado interno.
Un ejemplo clásico es la gestión de cuentas bancarias con protección contra sobregiros.

\begin{lstlisting}[caption={Clase BankAccount con validación de saldo.}, label={list:bank_account}]
class BankAccount:
    def __init__(self, owner, balance=0.0):
        self.owner = owner
        self.balance = balance

    def deposit(self, amount):
        if amount > 0:
            self.balance += amount
            print(f"Deposito de {amount} exitoso. Saldo actual: {self.balance}")
        else:
            print("El monto a depositar debe ser positivo.")

    def withdraw(self, amount):
        if amount > self.balance:
            print("Fondos insuficientes. Operacion cancelada.")
        elif amount <= 0:
            print("Monto de retiro invalido.")
        else:
            self.balance -= amount
            print(f"Retiro de {amount} exitoso. Saldo restante: {self.balance}")

cuenta = BankAccount("Juan", 1000)
cuenta.deposit(500)
cuenta.withdraw(2000)  # Muestra mensaje de fondos insuficientes
\end{lstlisting}

\subsection{Gestión del Estado Interno y Alternancia de Estados}
Los objetos pueden cambiar de estado mediante métodos dedicados a modificar variables booleanas o banderas internas.

\begin{lstlisting}[caption={Clase Light con métodos para alternar su estado.}, label={list:light}]
class Light:
    def __init__(self):
        self.is_on = False

    def turn_on(self):
        self.is_on = True

    def turn_off(self):
        self.is_on = False

    def toggle(self):
        self.is_on = not self.is_on
\end{lstlisting}

\section{Atributos con Estructuras de Datos Complejas}
Los atributos de una instancia no se limitan a tipos primarios (enteros, cadenas, flotantes); también pueden
almacenar colecciones completas como listas, diccionarios o conjuntos.

\begin{lstlisting}[caption={Clase Student almacenando una lista de notas.}, label={list:student}]
class Student:
    def __init__(self, name, marks):
        self.name = name
        self.marks = marks  # Lista de numeros flotantes o enteros

    def average(self):
        if not self.marks:
            return 0.0
        return sum(self.marks) / len(self.marks)

s1 = Student("Alice", [85, 90, 78, 92, 88])
print(f"Promedio de {s1.name}: {s1.average():.2f}")
\end{lstlisting}
